\documentclass[8pt]{beamer}

\setbeamertemplate{bibliography item}{\insertbiblabel}

%\usepackage{etex}
%\usecolortheme[RGB={205,173,0}]{structure} 
\usefonttheme{serif} 

%Pausing toggle
\def\mypause{
%\pause
} %Remove "%" or comment out

\usepackage{/Users/wdlinch3/Dropbox/Rashoumon/LaTeX/macros/WDL3beamer}

\setbeamertemplate{bibliography entry title}{}
\setbeamertemplate{bibliography entry location}{}
\setbeamertemplate{bibliography entry note}{}

%\usepackage{beamerthemesplit}

%%%%%%%%%%%%%%%%%%%%%%%%%%%%%%%%%%%%%%%%%%%%%%%%%%%
%%%%%%%%%%%%%%%%%%%%%%%%%%%%%%%%%%%%%%%%%%%%%%%%%%%
\title{Off-shell Supersymmetry and the M-theory Effective Action}
%\author{William D. Linch III \\
%~\\
%{\em George P. and Cynthia W. Mitchell Institute for\\
%Fundamental Physics and Astronomy, \\
%Texas A\&{}M University.}
%}
\author{William D. Linch III} 
\date{
{\em George P. and Cynthia W. Mitchell Institute for\\
Fundamental Physics and Astronomy, \\
Texas A\&{}M University.}\\
~\\
MIT CTP 11 September, 2019\\
~\\
Based on work 
%\cite{Becker:2016xgv, %ATH
%Becker:2016rku,	%NATH
%Becker:2016edk,	%G2Potential
%Becker:2017njd, 	%All CS
%Becker:2017zwe,	%Linearized
%Becker:2018phr}	%Current
with\\
K.\ Becker, 
M.\ Becker, 
D.\ Butter, 
S.\ Guha, 
S.\ Randall, \\
D.\ Robbins,
and
A.\ Sengupta.

}

\begin{document}

\frame{\titlepage}

\section[Outline]{}
%\frame{\tableofcontents}

%%%%%%%%%%%%%%%%%%%%%%%%%%%%%%%%%%%%%%%%%%%%%%%%%%%
\section{Introduction}


%%%%%%%%%%%%%%%%%%%%%%%%%
\frame{
\frametitle{Open string effective action}
\begin{itemize}
\item 
	Fradkin and Tseytlin (1985): Leading low-energy effective action of open string theory is Born-Infeld. 
%
\mypause
%
	Tseytlin's attempt at open superstring gives a weird answer. 
%
\mypause
%
\item
	Cecotti and Ferrara (1986): 4D, $N=1$ supersymmetry implies 
\begin{align}
\mathcal L_{CF} = \int d^2\theta W^2 + \int d^4\theta \,\Psi W^2\bar W^2
\nonumber
\end{align}
	They found a $\Psi= \Psi(D^2W^2, \bar D^2 \bar W^2) $ that reproduced Tseytlin's result but also observed that there is another corresponding to Born-Infeld again.
	% \cite{Cecotti:1986gb}.
%
\mypause
%
\item
	Metsaev, Rahkmanov, and Tseytin realize the mistake (R boundary condition) and recalculate with NS to find that the low-energy effective action is Born-Infeld (1987).
%
\mypause
%
\item
	Bagger and Galperin demonstrate that, starting in 4D, $N=1$ superspace, a second non-linearly-realized supersymmetry uniquely fixes the action to be Born-Infeld (1996).
Schematically,
\begin{align}
F^2 + F^4 + \textrm{8 supercharges} \Rightarrow \textrm{Born-Infeld}
\nonumber\end{align}
\end{itemize}
}






%%%%%%%%%%%%%%%%%%%%%%%%%
\frame{
\frametitle{Closed string effective action?}
Can we carry out the analogous steps for closed strings/gravity?
\begin{itemize}
\item
	Cannot do Tseytin {\it et al.}\ calculation because world-sheet theory is interacting (quartic) even for constant Riemann tensor. 

%
\mypause
%	

\item 
	Unclear what the gravitational analog of Born-Infeld is. Constant field strength? No birefringence? No ghosts? Determinantal form? Higher dimensions? Supersymmetric? \dots ?

%
\mypause
%

\item 
	Analog of  BG result may hold: There are only two supersymmetric densities in flat space. They would be related by $N\geq 2$ supersymmetry. Maybe even fixed by $N$ large enough?
\end{itemize}
}






%%%%%%%%%%%%%%%%%%%%%%%%%
\frame{
\frametitle{M-theory effective action?}
\begin{itemize}
\item
 	11D SG has maximal D and $N$.
%
\mypause
%

\item
 	First higher-derivative correction is of order (Riemann)${}^4$.

%
\mypause
%

\item
Higher-derivative correction requires a higher-derivative correction to the supersymmetry transformation, which requires higher derivatives in the effective action, {\it ad infinitum}.\\
 	It is believed that in this case the analog of the Born-Infeld statement holds:
\begin{align}
R + R^4 + \textrm{11D, $N=1$ supersymmetry} \Rightarrow \textrm{Leading M-theory LEEA?}
 \nonumber\end{align}

%
\mypause
%

\item
	Supersymmetrization order-by-order in components is unfeasible. Off-shell supersymmetry might fix this, but supersymmetry is on shell if $\#Q > 9$ (Berkovits 1993: Non-associativity of $\mathbb O$).

%
\mypause
%

\item 
	Perhaps keeping some $0<N<1$ manifest in 11D is enough?
\end{itemize}
}





%%%%%%%%%%%%%%%%%%%%%%%%%
\frame{
\frametitle{Previous work on fractional superspaces}
\begin{itemize}
\item 
	Marcus, Sagnotti, and Siegel (1983): 10D SYM in a $\mathbb R^{10|4}$. 
	The vector multiplet in this 10D, $N=1/4$ superspace consists of a real field $V(x, \theta, \bar \theta ; y)$ familiar from 4D, $N=1$ and an $SU(3)$ triplet of chiral superfields $\Phi_i(x, \theta; y)$. 
%\begin{align}
%R + R^4 + \textrm{11D, $N=1$ supersymmetry} \Rightarrow \textrm{M-theory effective action?}
% \end{align}
Action $\sim \int d^{10}x \int d^2 \theta W^\alpha W_\alpha +$ Chern-Simons terms.

%
\mypause
%

\item
	Berkovits (1994): Hybrid string  in this superspace and open string field theory using the result of MSS (1995).

%
\mypause
%

\item	
	Seiberg and Witten (1994): %``Solution'' of 4D, $N=2$ $SU(2)$ super-Yang-Mills. 
	Starting with 4D, $N=2$  SYM in terms of $N=1$ superfields, there is a D-term integral of the K\"ahler function $K(\bar \Phi, \Phi)$ and an F-term integral of the gauge kinetic term determined by a holomorphic gauge kinetic function $f(\Phi)$. 
	The second non-manifest supersymmetry determines the former in terms of the latter: $K = -\frac i2 f' \bar \Phi +\hc$\\ 
	Imposing $N=4$ fixes $f(\Phi)$ uniquely.

%
\mypause
%

\item
	5D, $N=1/2$ linearized supergravity was constructed with Luty and Phillips (2002) and used to compute loop corrections to slepton masses in a phenomenological analog of heterotic M-theory (with also Buchbinder, Gates, Goh, and Ng).

%
\mypause
%

\item 
	Extended to 11D SG with group at Texas A\&M.\\ 
%	 \cite{Becker:2016xgv, %ATH
%Becker:2016rku,	%NATH
%Becker:2016edk,	%G2Potential
%Becker:2017njd, 	%All CS
%Becker:2017zwe,	%Linearized
%Becker:2018phr}.	%Current
%This is work with 
%K.\ Becker, 
%M.\ Becker, 
%D.\ Butter, 
%S.\ Guha, 
%S.\ Randall, 
%D.\ Robbins
%and
%A.\ Sengupta.
~\\
\underline{Goal:} ~~~~Higher-derivative corrections to the M-theory effective action.
\end{itemize}
}





%%%%%%%%%%%%%%%%%%%%%%%%%%%%%%%%%%%%%%%%%%%%%%%%%%
\section{11D, $N=1/8$ Superspace}
\frame{
\frametitle{Constructing 11D, $N=1/8$ Superspace}
We construct 11D supergravity in an off-shell superspace $\bm M\sim \bm X\times Y$ with $\bm X$ a 4D, $N=1$ supermanifold and $Y$ a bosonic 7D manifold.

%
\mypause
%


This proceeds in six steps:
%
\mypause
%

%\begin{description}
\begin{itemize}
\item[1] Decompose 11D SG fields 
\begin{align}
\begin{array}{cccclcc}
	e_{\bm m}{}^{\bm a} &\to &
		e_{m}{}^{a} ~,~ 
		\mathcal A_m^i ~,~ 
		{\color {red} g_{ij} }
&\with & a, m = 0,1,2,3
\\
	\psi_{\bm m}^{\bm \alpha}&\to &
		\psi_{m}^{\alpha} ~,~ 
		\psi_{m}^{\alpha i} ~,~ 
		\psi_{j}^{\alpha } ~,~ 
		\psi_{j}^{\alpha i} ~,~
&\and & i,j = 1,\dots, 7	
\\ 
	C_{\bm {mnp}} &\to &C_{mnp}~,~ C_{mn \, i} ~,~ C_{m \, ij}~,~ C_{ijk}
&\and & \alpha = 1,2	
\end{array}
\nonumber
\end{align}
\end{itemize}
%description}
}




%%%%%%%%%%%%%%%%%%%%%%%%%%%%%%%%%%%%%%%%%%%%%%%%%%
\frame{
\frametitle{Step 2 (of 6): Reinterpreting the Metric on $Y$}
%\begin{description}
\begin{itemize}
\item[2] 
	Rewrite the metric {\color{red} $g_{ij}$} on $Y$ in terms of a 3-form (Bryant 1987, Hitchin, Joyce 2000) 
	\begin{align}
	\sqrt{g}g_{ij} = \frac1{144}\epsilon^{k_1\cdots k_7} \varphi_{i k_1k_2} \varphi_{j k_3 k_4} \varphi_{k_5k_6k_7}
	\nonumber 
	\end{align}
%
\mypause
%
	\begin{itemize}
	\item 
		This 3-form is ``stable'' when $g_{ij}$ is invertible. A stable 3-form on a 7 manifold is a(n almost) $G_2$ structure. 
%
\mypause
%
	\item 
		The torsion of this $G_2$ structure is defined by $\nabla_l \varphi_{ijk} = T_l{}^m (\ast \varphi)_{ijkm}$.
%
\mypause
%
	\item 
		If $T_i^j =0$, the holonomy of $\nabla$ is contained in $G_2$. The $49 = 1+ 7+ 14+ 27$ components of this tensor make up the torsion forms 
		\begin{align}
		d\varphi = \tau_0 \ast \varphi + \tau_1 \wedge \varphi + \ast \tau_3
		~~~\and~~~
		d(\ast \varphi) = \tau_1 \wedge \ast \varphi + \tau_2 \wedge \varphi
		\nonumber\end{align}
%
\mypause
%
	\item 
		If, furthermore, $Y$ is compact and simply-connected, then $\mathrm{Hol}(\nabla) = G_2$.
%
\mypause
%
	\item 
		The functional $\int_Y \sqrt{g(\varphi)} = \int_Y \varphi \wedge \ast \varphi$ is called the Hitchin (2000) functional. When $d\varphi = 0$ its equation of motion is $d(\ast \varphi) =0$ $\Leftrightarrow$ $G_2$ holonomy metrics (``Topological M-theory'' Dijkgraaf, Gukov, Neitzke, Vafa 2004). 
	\end{itemize}
%\end{description}
\end{itemize}
}





%%%%%%%%%%%%%%%%%%%%%%%%%%%%%%%%%%%%%%%%%%%%%%%%%%
\frame{
\frametitle{Step 3 (of 6): Superfield Embedding}
%\begin{description}
\begin{itemize}

\item[3] 
	Embed the bosonic fields into 11D, $N=1/8$ superfields:
%
\mypause
%
	\begin{itemize}
	\item 
	$e_{m}{}^{a}$ and $\psi_{m}^{\alpha}$ in a real 4-vector-valued superfield $U^a(x, \theta , \bar \theta;y)$ familiar from (conformal) 4D, $N=1$ supergravity,
%
\mypause
%
	\item 
	$\psi_{m}^{\alpha i}$ into a spinor superfield $\Psi_{\alpha i} (x, \theta , \bar \theta;y)$ in the $\bar {\bm 7}$ of $GL(7)$,
%
\mypause
%
	\item 
	$\mathcal A_m^i $ into a real superfield $\mathcal V^i (x, \theta , \bar \theta;y)$ forming a vector on $Y$, 
%
\mypause
%
	\item
	$C_{mnp}$ into a real scalar superfield $X(x, \theta , \bar \theta;y)$ forming a 0-form on $Y$,
%
\mypause
%
	\item
	$C_{mn \, i}$ into a chiral spinor $\Sigma_{\alpha i} (x, \theta ;y)$ forming a 1-form on $Y$,
%
\mypause
%
	\item
	$C_{m \, ij}$ into 21 real superfields $V_{ij} (x, \theta , \bar \theta;y)$ forming a 2-form on $Y$,
%
\mypause
%
	\item
	$C_{ijk}+i {\color{red} \varphi_{ijk}}$ into 35 chiral superfields $\Phi_{ijk} (x, \theta ;y)$ forming a 3-form on $Y$.
\end{itemize}
All these fields are superconformal primaries of the 4D, $N=1$ superconformal algebra. 
%\end{description}
\end{itemize}
}







%%%%%%%%%%%%%%%%%%%%%%%%%%%%%%%%%%%%%%%%%%%%%%%%%%
\frame{
\frametitle{Step 4 (of 6): Chern-Simons Action}
%\begin{description}
\begin{itemize}

\item[4] 
	The last five form a ``non-abelian tensor hierarchy'' of superfields.\footnote{Mathematically, this is a {\it Diff}(7)-equivariant chain complex of differential forms on $Y$ with values in a complex of (sheaves of) representations of the 4D, $N=1$ super-Poincar\'e algebra.} Such things have a unique Chern-Simons-like invariant supersymmetrizing the term $\int C \wedge dC \wedge dC$:
%
\mypause
%
	\begin{align}
	-12 \kappa^2 \mathcal L_{CS} &= i\int d^2 \theta \mathcal E\, \left[ i \Phi_{[0,3]}\wedge  \mathbb G_{[4,4]}  + \Sigma^\alpha_{[2,1]} \wedge \mathbb W_{[2,6] \alpha }\right] 
	\cr&
	+\int d^4 \theta E \left[ V_{[1,2]} \wedge \mathbb H_{[3, 5]} - X_{[3,0]} \mathbb F_{[1,7]}\right]
	+\hc 
	\nonumber\end{align}
%
\mypause
%

	\begin{itemize}
	\item 
		$\Phi\wedge \mathbb G = G \Phi \wedge d \Phi + \cdots$ contains ``$G_2$ superpotential'' (Gukov 1999). 
%
\mypause
%
	\item 
		$\Phi\wedge \mathbb G = \frac12 \Phi \wedge W^\alpha \wedge W_\alpha +\cdots$ also contains the holomorphic gauge-kinetic function $f(\Phi)  = \Phi$ familiar from Seiberg-Witten theory.
	\end{itemize}

%
\mypause
%

These terms and the others (10 total) are required for gauge invariance. 

%
\mypause
%


Unique invariant containing F-term integrals:
$\nexists$ gauge-invariant superpotential or gauge-kinetic function.
%\end{description}
\end{itemize}
}








%%%%%%%%%%%%%%%%%%%%%%%%%%%%%%%%%%%%%%%%%%%%%%%%%%
\frame{
\frametitle{Steps 5 and 6: Complete Action}
%\begin{description}
\begin{itemize}

\item[5] 
	There is a natural lift of the Hitchin functional to 11D, $N=1/8$ superspace:
	\begin{align}
	S_H = -\frac 3{\kappa^2} \int d^7 y \int d^4x  \int d^4\theta E (\bar GG)^{1/3} \sqrt{g(F)}\, f(x)
	~~~,~~~ x = {H_i g^{ij}(F) H_j \over (\bar GG)^{2/3}}
	\nonumber\end{align} 
	Uniquely fixed by local supersymmetry on $\bm X$, diffeomorphism invariance on $Y$, tensor hierarchy gauge invariance, engineering dimensions, conformal weights, $U(1)_R$ weights, and $GL(7)$ representation.

%
\mypause
%

\item[6] 
	Introducing the gravitino superfield, we can extend the local supersymmetry on $\bm X$ to all of $\bm M$. This fixes $f(x)$ exactly: 
	
%
\mypause
%

	The combination ${\hat f} := f -2 x f'$ satisfies the quartic equation 
	\begin{align}
		\frac x4 {\hat f}^4 +  {\hat f}^3 -1 = 0
			~~~\and~~~
		f(0)= 1
		~~~\Rightarrow
		~~~f(x)	
	\nonumber\end{align}
%	Together with the boundary condition $ {f}(0)= 1$, we can solve for $ {f}(x)$ in closed form.
%\end{description}
\end{itemize}
}






%%%%%%%%%%%%%%%%%%%%%%%%%%%%%%%%%%%%%%%%%%%%%%%%%%
\frame{
\frametitle{Checks}
We have subjected this result to the following consistency checks:
\begin{itemize}
\item
	Linearizing around a flat background, projecting to components, and throwing out fermions, we recover the correct bosonic component action of eleven-dimensional supergravity. 

%
\mypause
%

\item 
	Taking $\bm M = \mathbf R^{4|4} \times Y$ for general $Y$, reducing to components, throwing out all fermions and those bosons with polarizations in 4D directions, we obtain the ``scalar potential'' 
	\begin{align}
%\label{E:ScalarPotential1}
\hspace{-8mm}V=
	\tfrac14d_X^2 
	+ 2 g_{ij} \bm d^i \bm d^j 		
	-  d_{\mathbf 7 ij} ^2 
		+ \tfrac12 d_{\mathbf {14} ij}^2
	-\tfrac1{18} |f_{\mathbf 1 ijk}+\tfrac{i}2d_XF_{ijk}|^2
		-\tfrac1{24} |f_{\mathbf 7 ijk}|^2
		+\tfrac1{24} |f_{\mathbf {27} ijk}|^2
\nonumber
\end{align}
D- and F-term supersymmetry conditions $\Leftrightarrow$ $G_2$-holonomy and zero flux.
%
\mypause
%
\item 	
	Substituting the on-shell values of the auxiliary fields, this becomes 
\begin{align}
%\label{E:ScalarPotential}
V(x,y)&= 
%	-\tfrac{42}{18}\cdot \tfrac 9{16}\tau_0^2 
%	- \left(6+\tfrac{9\cdot 24}{24}\right)\tau_1^2 
%	+\tfrac12\cdot\tfrac14 \tau_2^2 
%	+\tfrac1{24} \tau_3^2			
%\cr&
%	+\left(\tfrac{49}{16} - \tfrac{42}{18} \cdot \tfrac 9{16}\right) \sigma_0^2 
%	+\left(2\cdot 9 - \tfrac{9\cdot 24}{24} \right) \sigma_1^2  
%	+\tfrac1{24} \sigma_3^2		
%\cr&= 
	-\tfrac{21}{16} \tau_0^2 
	- 15\tau_1^2 
	+\tfrac18 \tau_2^2 
	+\tfrac1{24} \tau_3^2			
%
	+\tfrac{7}{4} \sigma_0^2 
	+ 9 \sigma_1^2  
	+\tfrac1{24} \sigma_3^2	
\nonumber\end{align}
$G_2$ identities for torsion forms and $C$-field analogs $\Rightarrow$
%
\mypause
%
\begin{align}
%\label{E:ScalarPotential}
V(x,y)&= 
	-\tfrac12 R(g)
		+\tfrac1{4\cdot 4!} F_{ijkl}^2 
\nonumber\end{align}
$=$ Einstein-Hilbert and Maxwell terms on $Y$ from {\em auxiliary field mechanism}.
\end{itemize}
}






%%%%%%%%%%%%%%%%%%%%%%%%%%%%%%%%%%%%%%%%%%%%%%%%%%
\section{Higher derivatives}
%%%%%%%%%%%%%%%%%%%%%%%%%
\frame{
\frametitle{Higher derivatives?}
Can we guess higher-derivative terms? Many are easy. Most interesting one:
\begin{itemize}
\item We want to covariantize the term $\int C \wedge X_8$. 

%
\mypause
%

\item It contains $\int C \wedge R\wedge R\wedge R\wedge R$. 

%
\mypause
%

\item Decomposing, it contains $\int C \wedge W_X\wedge W_X\wedge W_Y\wedge W_Y$ 
%in terms of the Weyl tensors of $X$ and $Y$. 
(Weyl tensors of $X$, $Y$). 

%
\mypause
%

\item 
%The Weyl tensor on $X$ is in the gravitino curl $W_{\alpha \beta \gamma}$ familiar from 4D, $N=1$ supergravity. 
$W_X$ $\Leftrightarrow$ gravitino curl $W_{\alpha \beta \gamma}$ from 4D, $N=1$ supergravity.
So term $\Leftrightarrow$ %we expect this to come from 
$\varepsilon^{ijklmnp} \int \Phi_{ijk} W^{\alpha \beta \gamma}W_{\alpha \beta \gamma} S_{lmnp}$ for some $S\sim (W_Y)^2$ . %constructed from the square of the Weyl tensor on $Y$.

%
\mypause
%

\item But the Riemann tensor on $Y$ is real, so $S$ cannot be constructed directly from it. This was also the case for the potential, so perhaps it is generated by some auxiliary field mechanism?

%
\mypause
%

\item But the Weyl tensor is in the $\bm {70'}$ representation of $G_2$, and there is no auxiliary field for this!
\end{itemize}
}






%%%%%%%%%%%%%%%%%%%%%%%%%%%%%%%%%%%%%%%%%%%%%%%%%%
\frame{
\frametitle{Strategies}
We are pursuing a few strategies to study this problem. 
\begin{itemize}
\item[1] Consider a simpler problem:
%
\mypause
%
	\begin{itemize}
	\item 
		Construct Born-Infeld in 10D, $N=1/4$ superspace. 

%
\mypause
%

	\item 
		Auxiliary fields in $SU(3)$ reps:
			\begin{itemize}
			\item $\bar{\bm 3}$ $\leftrightarrow$ chiral $\Phi_i$ $\leftrightarrow$ $(2,0)$ part $F_{ij}$ of the curvature of the connection in a Calabi-Yau compactification. 
			\item a singlet $\leftrightarrow$ vector field $V$ $\leftrightarrow$ trace of $(1,1)$ part $F_k{}^k$. 
			\end{itemize}
			F-term SUSY $\Leftrightarrow$ holomorphic connex. \\
			D-term SUSY $\Leftrightarrow$ Donaldson-Uhlenbeck-Yau equation. 
%
\mypause
%

	\item The analog of the $\bm {70'}$ above is the $\bm 8={6\choose 2} -3-3-1$ $\leftrightarrow$ $F_i{}^j-\frac13 \delta_i{}^j F_k{}^k$ (the rest).
	\end{itemize}
%
\mypause
%
	Solution to the $\bm 8$ problem in BI  $\stackrel? \longleftrightarrow$ solution to $\bm {70'}$ problem in M-theory.
%
\mypause
%
	
	Better yet: If we can show there is no solution to the $\bm 8$ problem, we can safely forget about the $\bm {70'}$ problem and go work on something else.
\end{itemize}
}
%%%%%%%%%%%%%%%%%%%%%%%%%
\frame{
\frametitle{Strategies}
\begin{itemize}
\item [2${}^{\mathrm{nd}}$] Strategy:
Shut up and calculate
	\begin{enumerate}
	\item Construct off-shell supergeometry
	\item Construct $R\wedge R$ super-4-form
	\item Square to get $R\wedge R\wedge R\wedge R$
	\end{enumerate}

%
\mypause
%

%\item 
	To see how this works, first do easier 5D, $N=1/2$ superspace: 
%	In that case, the linearized supergeometry was already known from previous work with 
%
\mypause
%
	\begin{enumerate}
	\item Construct non-linear 5D, $N=1/2$ superspace supergravity {\ala} 11D, $N=1/8$. (Linearized limit agrees 2003 result.) 
%
\mypause
%

	\item We then use the known supergeometry (with Gates and Phillips 2003) to construct $R\wedge R$ (\cf Hanaki, Ohashi, and Tachikawa 2006):
\begin{align}
\mathbb G &= W^{\alpha \beta \gamma} W_{\alpha \beta \gamma} 
\cr
\mathbb H_i &= -\tfrac14 D^\gamma F_i^{\prime \alpha \beta} W_{\alpha \beta \gamma}
	+\tfrac 18 F_i^{\prime \alpha \beta}\bar D^{\dt \alpha}  D_\beta G_{\un a}
	+\tfrac 18 \bar D^{\dt \alpha} F_i^{\prime \alpha \beta}D_\beta G_{\un a} 
\cr	&	+\tfrac 1{16} \left( D^\beta \bar D_{\dt \alpha} F'_{\alpha \beta i} 
	+\tfrac 14 D_\alpha \bar D^2 \bar \lambda_{\dt \alpha i} \right) G^{\un a}
	+\tfrac 1{32} D^2 \lambda_i^\alpha D_\alpha R
	+ \mathrm{h.c.}
\nonumber\end{align}
%This is the superspace analog of the 2006 component calculation by 

%
\mypause
%
	\item This step is trivial in 5D. 
	\end{enumerate}
Conveniently, this answer can be lifted to 11D, $N=1/8$ unchanged. 
\end{itemize}	

}







%%%%%%%%%%%%%%%%%%%%%%%%%%%%%%%%%%%%%%%%%%%%%%%%%%
\section{Conclusion}
%%%%%%%%%%%%%%%%%%%%%%%%%
\frame{
\frametitle{Summary}
\begin{itemize}
\item 
	We have embedded eleven-dimensional supergravity in an off-shell 11D, $N=1/8$ superspace. 
\item 
	We have constructed the linearized off-shell supergeometry for it by starting with the fields and their gauge and local supersymmetry transformations and constructing the Riemann tensor, gravitino curls, and super-4-form field strength. 
\item 
	We have partially constructed the (Riemann)${}^2$ invariants with the intention of using them to construct the leading $\sim R^4$ M-theory corrections.
\item 
	Along the way, we have revisited the 5D, $N=1/2$ analog of this problem and are currently studying the 10D Born-Infeld analog. 
\end{itemize}
}

%%%%%%%%%%%%%%%%%%%%%%%%%
\frame{
\frametitle{Outlook}
\begin{itemize}
\item 
	Success in the construction of even the most non-trivial part of the $\int C \wedge X_8$ term is not the end of the story: 
	\begin{enumerate}
	\item 
		Chern-Simons-like action is not invariant under all 11D super-diffeomorphisms (\cf 2-derivative action). 

%
\mypause
%

	\item 
		D-term integral that fixes this problem already determined by $f(x)$. The non-invariance patched up by shift $G\to G + R\wedge R$ of tensor hierarchy field strengths (Witten 1996).
		% (this is what that term was for in the 2-derivative case and it was unique.

%
\mypause
%

	\item
		But D-term integrand is non-polynomial $\Rightarrow$ non-polynomial dependence in Riemann.
	\end{enumerate}
%
\mypause
%
Outcome expected from general arguments:\\
{\center 
\hspace{15mm} 
Supersymmetrization of $R^4$ must be non-polynomial.}
\end{itemize}
}






%%%%%%%%%%%%%%%%%%%%%%%%%%%%%%%%%%%%%%%%%%%%%%%%%%
\frame{
\frametitle{Outlook}
%\begin{itemize}
This would be an 11D supergravity analog of Born-Infeld. If this is correct, it would settle many questions about the definition of the latter but also leave some open and suggest new ones:
	\begin{itemize}
		\item It is not constructed purely from the curvature. Some derivatives thereof will be generated by the superspace integration. These are needed for supersymmetry.
%		\item There are no ghosts. 
		\item What is the fate of the singular GR solutions? This should not be difficult to answer once the action is known.
		\item Can the bosonic part be put in determinantal form? Does it define a new characteristic class?
%
\mypause
%
		\item Is it the leading low-energy effective action of M-theory? This could be checked by comparing terms available in the literature from {\it e.g.}\ amplitude computations. 
%
\mypause
%
		\item Is there a 64-supercharge theory for which the M-theory result is the Goldstone? Witten's 12D theory and $E_8$ framing (1996)? 
		\item $\cdots$
	\end{itemize}
%\end{itemize}
In addition, there are many applications/extensions (superstrings, compactifications, SFT, ...).

}






%%%%%%%%%%%%%%%%%%%%%%%%%%%%%%%%%%%%%%%%%%%%%%%%%%
\frame{
\frametitle{}
\begin{center}
{\bf \huge THANK YOU!}
\end{center}
}


%%%%%%%%%%%%%%%%%%%%%%%%%%%%%%%%%%%%%%%%%%%%%%%%%%
\frame{
\frametitle{NATH fields and transformations}

\begin{align}
\delta_0 \Psi_{\alpha i}
	&= \Xi_{\alpha i}
	+ \bar G\, \mathcal G_{i j} \nabla_\alpha \Omega^j
	+ i \,W_{\alpha i j} \,\bar{\Omega}^{j}
~,
\cr
\Delta_1 \Phi &= \Xi^\alpha \wedge W_\alpha
	+ \bar\nabla^2 \Big(
	\frac{i}{2} \iota_{\bar\Omega} \mathcal U
	+ \frac{1}{4} \iota_{\bar\Omega} (F \wedge H)
	\Big)
	~,\cr
% 	
\Delta_1 V &= - \frac{i}{2} \bar G \iota_\Omega F + \frac{i}{2} G \iota_{\bar\Omega} F ~,
\cr
% 
\Delta_1 \Sigma_\alpha &= -G\, \Xi_{\alpha}~, 
\cr
% 
\Delta_1 X &= -\frac{i}{2} \bar G\, \iota_\Omega H + \frac{i}{2} G\, \iota_{\bar\Omega} H~,
\cr
% 
\delta_1 \mathcal V &= -\frac{1}{2} \bar G\, \Omega - \frac{1}{2} G\, \bar \Omega ~.
\nonumber
\end{align}
\begin{align}
\mathcal U_{ijkl} &:= 3 \,\epsilon_{ijklmnp} \frac{\partial}{\partial F_{mnp}} \Big( \sqrt g\,  {f}\, (G \bar G)^{1/3}\Big) 
%\cr
%	&= \frac{1}{3!} \sqrt{g}\,\epsilon_{ijklmnp} \Big(
%	F^{mnp} (\cF + |H|^2 \cF')
%	- 9 \cF'\, H^{m} F^{n p q} H_q (G \bar G)^{-2/3}
%	\Big) (G \bar G)^{1/3}~.
\cr
\mathcal G_{i j} 
	&= (G \bar G)^{-1/3} \Big(\hat  {f}^{-1}g_{i j} + \frac i2 (G \bar G)^{-1/3} \hat  {f} F_{i j k} \,g^{k l} \,H_l \Big)
\nonumber
\end{align}
\begin{align}
\delta_0 \Psi_{\alpha i} &= \Xi_{\alpha i} + g_{i j} D_\alpha \Omega^j  ~,
\cr
\delta_1 \Phi_{ijk} &= \tfrac1{2i} \tilde \varphi_{ijkl} \bar D^2 \bar \Omega^l ~,
\cr
\delta_1  V_{ij} &= \tfrac1{2i} \varphi_{ijk} (\Omega^k -\bar \Omega^k) ~,
\cr
\delta_1  \Sigma_{\alpha i}&= -\Xi_{\alpha i} ~,
\cr
\delta_1 {\mathcal V}{}^i &= -\tfrac12(\Omega^i + \bar \Omega^i)~.
\nonumber
\end{align}
}

%%%%%%%%%%%%%%%%%%%%%%%%%%%%%%%%%%%%%%%%%%%%%%%%%%
\frame{
\frametitle{NATH field strengths and Bianchi identities}
\begin{align}
E &= \partial \Phi
	\cr
F &= \tfrac1{2i}\left( \Phi - \bar \Phi\right) - \partial V
	\cr
W_\alpha &= -\tfrac14 \bar \nabla^2 \nabla_\alpha V 
	+\partial \Sigma_\alpha 
	+ \iota_{\mathcal W_\alpha} \Phi
	\cr
H &= \tfrac1{2i}\left(\nabla^\alpha \Sigma_\alpha - \bar \nabla_{\dt \alpha} \bar \Sigma^{\dt \alpha} \right)
	-\partial X 
	-\omega_{\mathsf h}(\mathcal W, V)
	\cr
G&= -\tfrac14 \bar \nabla^2 X 
	+ \iota_{\mathcal W^\alpha} \Sigma_\alpha	
\nonumber
\end{align}
\begin{align}
0 &= - \partial E
	\cr
\tfrac1{2i}\left(E - \bar E\right) &=  \partial F
	\cr
-\tfrac14 \bar {\nabla}^2 {\nabla}_\alpha F  &= 
	-\partial W_\alpha 
	- \iota_{\mathcal W_\alpha} E
	\cr
\tfrac1{2i}\left({\nabla}^\alpha W_\alpha - \bar {\nabla}_{\dt \alpha} \bar W^{\dt \alpha} \right)&= 
	\partial H 
	+ \omega_{\mathsf h}(\mathcal W, F)
	\cr
-\tfrac14 \bar {\nabla}^2 H &= 
	-\partial G 
	- \iota_{\mathcal W^\alpha} W_\alpha 
	\cr
\bar {\nabla}_{\dt \alpha} G &= 0
\nonumber
\end{align}
\begin{align}
	\omega_{\mathsf h}(\chi_\alpha, v) := 
\iota_{\chi^\alpha} {\nabla}_\alpha v 
	+ \iota_{\bar \chi_{\dt \alpha}} \bar {\nabla}^{\dt \alpha} v 
	+\tfrac12 \left(  \iota_{{\nabla}^\alpha\chi_\alpha} v 
		+ \iota_{\bar {\nabla}_{\dt \alpha} \bar \chi^{\dt \alpha}} v 
	\right)
\nonumber
\end{align}	

}

%%%%%%%%%%%%%%%%%%%%%%%%%%%%%%%%%%%%%%%%%%%%%%%%%%
\frame{
\frametitle{Chern-Simons-like invariant}
\begin{align}
-12 L_{CS} &= 
i \int d^2 \theta  \, \mathcal E\,   \Phi \left[
	E G + \tfrac12 W^\alpha W_\alpha - \tfrac i4 \bar {\nabla}^2 (  F H) 
	\right]
\cr&+ \int d^4 \theta  \, E\,  V \left[  
	\left( E + \bar E\right) H + \omega(W, F) 
	- i {\nabla}^\alpha F\iota_{\mathcal W_\alpha}F 
	+ i \bar {\nabla}_{\dt \alpha} F \iota_{\bar{\mathcal W}^{\dt \alpha} }F 
	\right]
\cr&+i \int d^2 \theta  \, \mathcal E\, \Sigma^\alpha \left[
	E W_\alpha -\tfrac i4 \bar {\nabla}^2 (  F {\nabla}_\alpha F) 
	\right]
\cr&- \int d^4 \theta  \, E\, X \left[  
	\left( E + \bar E\right) F
	\right]
	+\mathrm{h.c.} 	
\nonumber
\end{align}
}

%%%%%%%%%%%%%%%%%%%%%%%%%%%%%%%%%%%%%%%%%%%%%%%%%%%
%\frame{
%\frametitle{Born-Infeld}
%
%
%
%}





%%%%%%%%%%%%%%%%%%%%%%%%%%%%%%%%%%%%%%%%%%%%%%%%%%
{\tiny
%\small
%\footnotesize
%\linespread{1.1}\selectfont
\bibliography{/Users/wdlinch3/Dropbox/Rashoumon/LaTeX/BibTex/BibTex}
\bibliographystyle{unsrt}
} % end resize


\end{document}